\section{Role-specific framework and runtime inventory}\label{sec:C22:inventory}
Table~\ref{tab:C22:inventory} assigns each listed framework or managed runtime a concrete inspection target. The entries describe documentation surfaces, not measured capability or comparative quality. The final column is a discriminating test: it identifies an observation needed before an engineering conclusion is justified. Where a source redirects, the resolved resource defines the inspected family; historical names are lineage identifiers. Package and hosted-service parity remain unverified unless established by a version-specific test.
\begingroup\setlength{\tabcolsep}{3pt}\small
\bookneedspace{.42\textheight}
\begin{longtable}{L{25mm}L{35mm}L{48mm}}
\caption{Framework and managed-runtime surfaces with discriminating checks.}\label{tab:C22:inventory}\\
\toprule Product & Documented surface & Required test \\
\midrule\endfirsthead
\toprule Product & Documented surface & Required test \\
\midrule\endhead
LangGraph\cite{executionproduct001} & Graph checkpoints and stores & Crash after an effect; distinguish restored state from repeated dispatch. \\
OpenAI Agents SDK\cite{executionproduct002} & Agent and Runner loop interfaces & Preserve rejected calls, handoff ownership and application side-effect receipts. \\
Agent Development Kit — ADK\cite{executionproduct003} & Workflow and task delegation objects & Compare graph transitions under retry, human interruption and version mismatch. \\
PydanticAI\cite{executionproduct004} & Typed loop and configurable harness & Inject validly typed but unauthorized arguments; separate types from permission checks. \\
Deep Agents\cite{executionproduct006} & Filesystem offload and isolated subagents & Remove a rare constraint during compaction and test later recovery. \\
Microsoft Agent Framework\cite{executionproduct007} & Multi-language agents and workflows & Round-trip state across language adapters; retain cancellation and error semantics. \\
Agno / AgentOS\cite{executionproduct010} & Agent teams and runtime interfaces & Test whether shared memory commits conflict or silently overwrite one another. \\
LlamaAgents / Agent Workflows\cite{executionproduct011} & Event-driven workflow composition & Duplicate an input event and inspect whether downstream effects repeat. \\
Strands Agents / Harness SDK\cite{executionproduct012} & Composable harness control interfaces & Locate the actual agent-loop implementation after the source redirect; pin the adapter. \\
OpenClaw\cite{executionproduct014} & Cross-channel agent gateway & Revoke one channel principal and test isolation of historical context. \\
Mastra\cite{executionproduct019} & TypeScript agents and workflows & Interrupt a suspended workflow and verify that resumption preserves evidence identity. \\
Claude Agent SDK\cite{executionproduct020} & Claude-oriented SDK execution surfaces & Identify application versus provider ownership of hooks, session state and tool effects. \\
Haystack\cite{executionproduct026} & Retrieval pipelines and agent components & Replace retrieval with oracle evidence while holding generation and tool budgets fixed. \\
NeMo Fabric\cite{executionproduct031} & Cross-harness execution and telemetry & Check whether adapter traces preserve truncation, delegated work and error provenance. \\
smolagents\cite{executionproduct032} & Code-action agent loop & Inspect the generated-code execution boundary and sandbox mediation of effects. \\
DeepSeek Harness\cite{executionproduct035} & Plugin-oriented harness interfaces & Disable one plugin and inspect transitive tool exposure and persisted assumptions. \\
AgentScope\cite{executionproduct039} & Agent application and team abstractions & Test cancellation propagation and whether conversation identifiers imply durable storage. \\
CrewAI\cite{executionproduct041} & Role-based crews and explicit flows & Compare identical role assignments with deterministic flow and model-selected delegation. \\
Pi\cite{executionproduct049} & Terminal harness and extension interfaces & Load a minimal extension and audit which credentials and context it receives. \\
DSPy\cite{executionproduct052} & Program optimization against evaluation & Keep an untouched holdout and count selection attempts as part of optimization cost. \\
OGX, formerly Llama Stack\cite{executionproduct054} & Portable server-side agent interfaces & Follow the OGX lineage and verify old client schema compatibility explicitly. \\
Qwen-Agent\cite{executionproduct056} & Qwen-oriented tools and retrieval & Match tool serialization and model settings before attributing benchmark differences. \\
n8n\cite{executionproduct057} & Workflow automation with agent steps & Separate durable workflow state from model-visible memory and connector permissions. \\
Langflow\cite{executionproduct062} & Visual Python agent composition & Export the executable graph and test whether visual and API paths agree. \\
Dify\cite{executionproduct066} & Visual workflows and retrieval pipelines & Check which runtime and connector components are actually inspectable and licensed. \\
Elastic Agent Builder\cite{executionproduct068} & Search-grounded agent construction & Change document permissions after indexing and test retrieval-time enforcement. \\
CAMEL\cite{executionproduct071} & Collaborating agent framework & Compare communicating workers against equal-cost independent samples and one-agent controls. \\
Coze Studio\cite{executionproduct082} & Self-hostable visual agent workflows & Inspect credential scoping, deployment dependencies and duplicate-event recovery. \\
UiPath Maestro\cite{executionproduct018} & Business-process orchestration patterns & Interrupt a human approval chain; bind the resumed action to current authorization. \\
Agent Bricks\cite{executionproduct042} & Custom-agent hosting and data integration & Verify workspace persistence and trace identity across a container restart. \\
Bedrock AgentCore\cite{executionproduct043} & Composable runtime identity and memory services & Test service boundaries separately; a memory feature does not establish agent correctness. \\
Gemini Enterprise Agent Platform\cite{executionproduct044} & Managed agent lifecycle services & Distinguish low-code, managed API and custom ADK deployment semantics. \\
Agents API\cite{executionproduct060} & Managed sessions and execution harness & Compare saved turn history with actual external-effect reconciliation under failure. \\
Claude Managed Agents\cite{executionproduct061} & Managed long-running agent sessions & Inspect environment persistence, permission policy and incomplete-turn outcomes independently. \\
Agents \& Conversations API\cite{executionproduct089} & Agent configuration and conversation entries & Test history continuity separately from tool execution and external transaction guarantees. \\
\bottomrule\end{longtable}
\endgroup


Several entries cross category boundaries. Pi can be used as a coding harness and as an extension substrate. NeMo Fabric concerns execution and telemetry across harnesses rather than a replacement for the model policy. DSPy concerns program and prompt optimization against evaluation objectives; it should be tested for selection overfitting as well as task quality. A visual builder such as Langflow, Dify, Coze Studio, or n8n can host model-driven steps inside deterministic workflows. Its graphical interface does not remove the need to examine stored state, credential scope, and replay behavior.

Managed platforms change the observation boundary. AgentCore, Agent Bricks, and Google's enterprise platform document infrastructure surfaces that an application can compose with custom agents. Managed agent APIs supply a larger portion of the loop. These roles should be compared through migration effort, observability, failure recovery, and enforceable policy controls, rather than counted as additional reasoning architectures. A serverless runtime's availability commitment is not evidence that an agent's task output is correct.

